\chapter{Digital Implementation and Practical Considerations}
\label{chap:digital}

\begin{learningobjective}
\begin{itemize}[leftmargin=*]
    \item Understand how sampling, quantization, and computational delays affect closed-loop performance.
    \item Translate continuous-time controllers into numerically stable discrete implementations.
    \item Plan validation workflows from simulation to hardware-in-the-loop testing.
\end{itemize}
\end{learningobjective}

\section{Concept Map for Beginners}
\begin{conceptroadmap}
\begin{enumerate}[leftmargin=*]
    \item Look at sampling as the act of taking snapshots; relate Nyquist and aliasing to everyday examples like misinterpreted wagon wheels.
    \item Translate controllers carefully so that the discrete implementation preserves the intent of the continuous design.
    \item Build confidence stepwise—simulate, run on the processor, then integrate hardware—so each stage validates assumptions about timing and quantization.
\end{enumerate}
\end{conceptroadmap}

\section{Concept--Math--Engineering Loop}
\begin{table}[h]
    \centering
    \small
    \renewcommand{\arraystretch}{1.2}
    \begin{tabularx}{\linewidth}{|>{\raggedright\arraybackslash}p{0.28\linewidth}|>{\raggedright\arraybackslash}p{0.34\linewidth}|>{\raggedright\arraybackslash}X|}
        \hline
        \textbf{Concept} & \textbf{Mathematics} & \textbf{Engineering Practice} \\
        \hline
        Sampling effects & Nyquist limit, discrete pole mapping $z = e^{\lambda T_s}$ & Reproduce aliasing plots with \tutorialcmd{digital}. \\
        \hline
        Discretization fidelity & Euler and Tustin substitutions & Compare discrete PID implementations by switching methods in the tutorial. \\
        \hline
        Nonidealities & Quantization models, delay-induced phase lag & Use simulator options for quantization/jitter and inspect generated figures for margin loss. \\
        \hline
    \end{tabularx}
    \caption{Connecting digital-control concepts to the math and experimentation workflow.}
\end{table}

\section{Sampling Theory Refresher}
\begin{beginnerguide}
    \item[\textbf{What?}] Reviews how sampling, aliasing, and pole mapping affect discrete-time representations.
    \item[\textbf{Why?}] Sampling choices directly impact closed-loop stability and fidelity when migrating to digital controllers.
    \item[\textbf{When?}] Before discretising controllers or choosing hardware update rates.
    \item[\textbf{How?}] Apply Nyquist limits, map continuous poles via $z=e^{\lambda T_s}$, and simulate aliasing scenarios.
    \item[\textbf{Example}] Use the tutorial to sample a 20~Hz plant at 25~Hz and observe the resulting 5~Hz alias in the output spectrum.
\end{beginnerguide}
Sampling a continuous signal $x(t)$ at period $T_s$ yields $x[k] = x(k T_s)$. Key implications:
\begin{itemize}[leftmargin=*]
    \item Nyquist criterion requires sampling frequency $\omega_s = 2 \pi / T_s$ to exceed twice the highest significant frequency component ($\omega_s > 2 \omega_{\max}$).
    \item Continuous-time poles $\lambda_i$ map to discrete poles $z_i = \mathrm{e}^{\lambda_i T_s}$. Poles with large negative real parts may cluster near the origin if $T_s$ is large, distorting dynamics.
    \item Input/output delays $\tau$ translate to factors of $z^{-d}$ with $d = \lfloor \tau / T_s \rceil$, reducing phase margin.
\end{itemize}
\subsection{Worked Example: Aliasing and Timing Jitter}
Consider a plant whose dominant mode oscillates at $20$~Hz ($\omega_{\max} = 40 \pi$~rad/s). Sampling at $T_s = 0.04$~s (25~Hz) violates Nyquist because $\omega_s = 50 \pi < 2 \omega_{\max}$. The aliased frequency observed by the controller is
\[
    f_\text{alias} = |f_\text{signal} - N f_s| = |20 - 1 \cdot 25| = 5~\text{Hz},
\]
causing the loop to chase a much slower oscillation and mis-tune gains. Edit \texttt{tutorial\_digital} to set the sampling frequency to $25$~Hz and re-run
\begin{lstlisting}[language=bash]
python -m scripts.tutorials digital --output-dir figures
\end{lstlisting}
to regenerate \texttt{figures/digital\_sampling\_spectrum.png} with both the original 20~Hz peak and the aliased 5~Hz component.

Execution jitter produces an equivalent phase lag. If a task scheduled for $T_s = 10$~ms experiences a worst-case delay of $\delta = 2$~ms, the additional phase lag near crossover frequency $\omega_c$ is approximately $\Delta \phi \approx -\omega_c \delta$. For $\omega_c = 30$~rad/s, the lag reaches $-3.4^{\circ}$. Extend \texttt{tutorial\_digital} to draw a random delay (e.g., uniform in $[0, \delta]$) before each controller evaluation and log the resulting margins so that \texttt{figures/digital\_jitter\_margin.png} captures the degradation.

\section{Controller Discretization}
\begin{beginnerguide}
    \item[\textbf{What?}] Methods for converting continuous controllers into stable, accurate discrete implementations.
    \item[\textbf{Why?}] Naïve discretisation can destabilise loops; understanding mapping choices preserves design intent.
    \item[\textbf{When?}] After designing continuous compensators and before coding them in embedded firmware.
    \item[\textbf{How?}] Choose an $s\mapsto z$ mapping (forward/backward Euler, Tustin), derive discrete transfer functions, and test numerically.
    \item[\textbf{Example}] Discretise a lead compensator with Tustin and compare with the exact ZOH discretisation produced by SciPy’s \texttt{cont2discrete}.
\end{beginnerguide}
Common $s \mapsto z$ mappings:
\begin{align}
    s &\approx \frac{z - 1}{T_s} && \text{(forward Euler)},\\
    s &\approx \frac{z - 1}{z T_s} && \text{(backward Euler)},\\
    s &\approx \frac{2}{T_s} \frac{z - 1}{z + 1} && \text{(Tustin/bilinear)}.
\end{align}
Forward Euler is simple but may destabilize high-frequency dynamics; Tustin preserves stability and usually offers better fidelity. Always implement derivative terms with low-pass filtering to avoid amplifying noise.

\section{Quantization, Saturation, and Noise}
\begin{beginnerguide}
    \item[\textbf{What?}] Practical nonidealities introduced by finite resolution and actuator limits.
    \item[\textbf{Why?}] These effects can erode performance or cause instability if ignored during design.
    \item[\textbf{When?}] During controller implementation and hardware testing.
    \item[\textbf{How?}] Model quantisation as rounding, include saturation blocks, and simulate noise to evaluate robustness.
    \item[\textbf{Example}] Run the simulator with \texttt{quantization=1e-4} and actuator limits to see how PID integral windup appears without anti-windup logic.
\end{beginnerguide}
Finite resolution introduces quantization noise. Simulate the effects with the harness:
\begin{lstlisting}[language=python, caption={Quantization impact}]
result = simulate_discrete_system(
    Ad, Bd, C,
    controller_step,
    ts=Ts,
    steps=3000,
    reference=lambda k: 0.05 if k > 50 else 0.0,
    measurement_noise_std=5e-4,
    quantization=1e-4,
    actuator_limits=(-5.0, 5.0),
)
print(result.performance_metrics())
\end{lstlisting}
Integral action must incorporate anti-windup logic to prevent integrator windup when actuators saturate. Options include integral state clamping or back-calculation.

\subsection{Code Mapping}
\begin{table}[h]
    \centering
    \small
    \renewcommand{\arraystretch}{1.2}
    \begin{tabularx}{\linewidth}{|>{\raggedright\arraybackslash}p{4cm}|X|X|}
        \hline
        \textbf{Python} & \textbf{Formula} & \textbf{Explanation} \\
        \hline
        Simulator option & $y_k^\text{meas} = \text{round}(y_k/\Delta) \Delta$ & Uses the quantization flag in \texttt{simulate\_discrete\_system} to enforce the step size $\Delta$. \\
        \hline
        \texttt{PIDController.step()} & Anti-windup logic & Includes saturation checks and integral clamping/back-calculation. \\
        \hline
    \end{tabularx}
    \caption{Digital implementation effects mapped to simulation code.}
\end{table}

\section{Tutorial: Digital Control Experimentation}
\begin{beginnerguide}
    \item[\textbf{What?}] Scripted experiments exploring sampling, delay, and quantisation effects on PID loops.
    \item[\textbf{Why?}] Hands-on trials highlight the performance penalties of implementation constraints.
    \item[\textbf{When?}] Before deploying code to hardware or when documenting digital design trade-offs.
    \item[\textbf{How?}] Execute \tutorialcmd[--output-dir figures]{digital}\\and compare ideal versus nonideal scenarios.
    \item[\textbf{Example}] Generate \texttt{figures/digital\_pid\_responses.png} to compare nominal and delayed executions, noting increased oscillations due to jitter.
\end{beginnerguide}
\begin{enumerate}[leftmargin=*]
    \item Generate the digital-control figures with\\
    \tutorialcmd[--output-dir figures]{digital}.
    \item Study Figure~\ref{fig:digital-responses} to compare ideal, quantized, and delayed PID loops; note how quantization amplifies steady-state error while delay introduces oscillation.
    \item Examine Figure~\ref{fig:digital-spectrum} to see how different sampling periods reshape the output spectrum, echoing the sampling guidelines of \textcite{astrom2006advanced}.
    \item Inspect Figure~\ref{fig:digital-controls} to assess actuator usage under each scenario and motivate anti-windup strategies.
\end{enumerate}
Command-line recipes for rebuilding these figures and compiling the notes are summarised in Appendix~\ref{chap:tooling}.
\IfFileExists{../figures/digital_pid_responses.png}{
    \begin{figure}[h]
        \centering
        \includegraphics[width=0.75\linewidth]{../figures/digital_pid_responses.png}
        \caption{Impact of quantization and sensor delay on PID tracking.}
        \label{fig:digital-responses}
    \end{figure}
}{
    \begin{figure}[h]
        \centering
        \fbox{\parbox{0.7\linewidth}{Run \tutorialcmd{digital}\\to create \texttt{figures/digital\_pid\_responses.png}, then recompile.}}
        \caption{Placeholder for the digital PID response comparison.}
        \label{fig:digital-responses}
    \end{figure}
}
\IfFileExists{../figures/digital_pid_controls.png}{
    \begin{figure}[h]
        \centering
        \includegraphics[width=0.75\linewidth]{../figures/digital_pid_controls.png}
        \caption{Control effort under ideal, quantized, and delayed execution.}
        \label{fig:digital-controls}
    \end{figure}
}{
    \begin{figure}[h]
        \centering
        \fbox{\parbox{0.7\linewidth}{Generate \texttt{figures/digital\_pid\_controls.png} via the digital tutorial to visualise actuator demands.}}
        \caption{Placeholder for the digital control-effort comparison.}
        \label{fig:digital-controls}
    \end{figure}
}
\IfFileExists{../figures/digital_sampling_spectrum.png}{
    \begin{figure}[h]
        \centering
        \includegraphics[width=0.75\linewidth]{../figures/digital_sampling_spectrum.png}
        \caption{Output spectra for multiple sampling periods.}
        \label{fig:digital-spectrum}
    \end{figure}
}{
    \begin{figure}[h]
        \centering
        \fbox{\parbox{0.7\linewidth}{Generate \texttt{figures/digital\_sampling\_spectrum.png} via the digital tutorial to examine sampling effects.}}
        \caption{Placeholder for the sampling-spectrum comparison.}
        \label{fig:digital-spectrum}
    \end{figure}
}

\section{Scheduling and Real-Time Constraints}
\begin{beginnerguide}
    \item[\textbf{What?}] Considerations for meeting timing deadlines and mitigating jitter in real-time systems.
    \item[\textbf{Why?}] Missed deadlines or variable execution times translate into effective delays that can destabilise control loops.
    \item[\textbf{When?}] While allocating tasks to processors or designing RTOS schedules.
    \item[\textbf{How?}] Profile computation, budget slack, add watchdogs, and design with extra phase margin to tolerate jitter.
    \item[\textbf{Example}] Estimate phase margin loss using $\Delta\phi \approx -\omega_c \delta$ for a task with \SI{2}{ms} jitter at $\omega_c=30$~rad/s.
\end{beginnerguide}
Digital controllers execute on embedded processors with finite resources:
\begin{itemize}[leftmargin=*]
    \item \textbf{Computation time}: the control algorithm must finish within $T_s$. Profiling and fixed-point optimization may be necessary on microcontrollers.
    \item \textbf{Jitter}: Variations in execution time act as multiplicative uncertainty. Design with additional phase margin to accommodate jitter.
    \item \textbf{Latency}: Sensor-to-actuator delays reduce effective phase margin; compensate via Smith predictors or redesign with lower bandwidth.
\end{itemize}

\section{Validation Pipeline}
\begin{beginnerguide}
    \item[\textbf{What?}] A staged approach (SIL $\rightarrow$ PIL $\rightarrow$ HIL) for verifying controllers before deployment.
    \item[\textbf{Why?}] Structured validation catches issues early, reducing risk during hardware trials.
    \item[\textbf{When?}] After implementing the controller in software and before integration with the physical plant.
    \item[\textbf{How?}] Simulate with software-in-the-loop, execute on target hardware with plant models, then connect to real sensors/actuators.
    \item[\textbf{Example}] Run SIL simulations with quantisation enabled, then execute the same control law on a development board in PIL mode to compare timing.
\end{beginnerguide}
\begin{enumerate}[leftmargin=*]
    \item \textbf{Software-in-the-loop (SIL)}: run Python simulations that include quantization, saturation, noise, and network-related effects.
    \item \textbf{Processor-in-the-loop (PIL)}: execute compiled control code on the target processor while interfacing with a simulator to measure timing and numerical precision.
    \item \textbf{Hardware-in-the-loop (HIL)}: connect actual hardware sensors and actuators (or a real-time plant emulator) to validate end-to-end behavior before deploying to the real system.
\end{enumerate}
Document the validation steps in the appendix, embedding diagnostic plots and logging procedures.

\section{Implementation Checklist}
\begin{beginnerguide}
    \item[\textbf{What?}] A quick audit of numerical and safety considerations prior to release.
    \item[\textbf{Why?}] Ensures critical safeguards—overflow protection, watchdogs, reproducible builds—are in place.
    \item[\textbf{When?}] Immediately before deploying firmware or handing off to quality assurance.
    \item[\textbf{How?}] Walk through the checklist items: scaling, anti-windup, watchdog timers, configuration management.
    \item[\textbf{Example}] Confirm integral states are clamped to actuator limits and that configuration files record the exact PID gains shipped.
\end{beginnerguide}
\begin{itemize}[leftmargin=*]
    \item Verify numerical robustness (overflow, quantization, scaling) for fixed-point implementations.
    \item Use watchdog timers and safety interlocks to handle communication failures or saturation events.
    \item Maintain reproducible build scripts and configuration management so that controller parameters match the documentation.
\end{itemize}

\begin{takeaway}
\begin{itemize}[leftmargin=*]
    \item Digital implementation details significantly influence real-world performance; treat them as part of the control design process.
    \item Proper discretization, anti-windup, and scheduling analysis prevent instability caused by sampling and computational effects.
    \item A staged validation strategy (SIL $\rightarrow$ PIL $\rightarrow$ HIL) mitigates risk before deploying controllers to physical hardware.
\end{itemize}
\end{takeaway}

\section{Module Review}
\begin{itemize}[leftmargin=*]
    \item \textbf{Sampling}: check Nyquist ($\omega_s > 2\omega_{\max}$) and preview aliasing via Figure~\ref{fig:digital-spectrum}; adjust $T_s$ accordingly.
    \item \textbf{Quantisation/delay}: evaluate tracking degradation (Figure~\ref{fig:digital-responses}) and control effort (Figure~\ref{fig:digital-controls}); implement anti-windup where saturation is visible.
    \item \textbf{Workflow}: follow the SIL $\rightarrow$ PIL $\rightarrow$ HIL path; capture logs and reproduce the tutorial to document scenarios.
    \item \textbf{Practice}: run \tutorialcmd{digital} and explore network-induced jitter or sensor noise to build intuition before hardware deployment.
    \item \textbf{Reading}: \textcite{astrom2006advanced} Ch.~11 for sampled-data issues and \textcite{ljung2012control} Ch.~12 for implementation case studies.
\end{itemize}

\section{Keyword Review}
\begin{vocabulary}
\begin{itemize}[leftmargin=*]
    \item \textbf{Sampling period $T_s$}: time between controller updates; sets discrete-time dynamics.
    \item \textbf{Aliasing}: misinterpretation of high-frequency signals when $T_s$ violates Nyquist.
    \item \textbf{Jitter}: variation in execution timing that behaves like additional delay.
    \item \textbf{Quantization}: rounding of signals to discrete levels, introducing effective noise.
    \item \textbf{Anti-windup}: strategy to prevent integrator accumulation during saturation.
    \item \textbf{HIL testing}: hardware-in-the-loop validation combining real controllers with simulated plants.
\end{itemize}
\end{vocabulary}
